% ============================================================================
%  chapter11.tex
%  Concurrency with Threads
% ============================================================================

\chapter{Concurrency with Threads}

\section{A Process with Many Simultaneous Selves}

Chapter~10 taught one process to serve many channels by never blocking.
This chapter teaches the complementary approach: give one process
\emph{many threads of execution}. A \emph{thread} is a single sequence
of instructions running inside a process. A process with three threads
is, from the CPU's point of view, three independent instruction
streams---each with its own stack and registers---that happen to share
one address space.

\begin{verbatim}
   Two views of the same program:

   PROCESSES (ch.5-9):               THREADS (this chapter):

   +----------+  +----------+       +-----------------------+
   | stack    |  | stack    |       | stack T1  (private)   |
   | heap     |  | heap     |       | stack T2  (private)   |
   | data     |  | data     |       | stack T3  (private)   |
   | text     |  | text     |       |   THE SHARED HEAP     |
   +----------+  +----------+       |   THE SHARED GLOBALS  |
        separate address             |   text                |
        spaces, no shared             +-----------------------+
        memory                         ONE address space,
                                       many instruction streams

   Threads share: the heap, the globals, the open descriptors.
   Threads do NOT share: their stacks, their registers.

   Sharing is the whole point -- and the whole danger.
\end{verbatim}

The trade-off with processes is now concrete. Processes are isolated
but must use pipes or memory mapping to share anything; threads share
everything trivially---a thread can read and write a global variable
the way any function can---but that sharing is exactly where races
come from. The rest of this chapter is the discipline that makes
sharing safe.

\begin{osconcept}
Threads are the kernel (or the thread library) multiplexing several
instruction streams onto one process. Because they share one address
space, the \emph{context switch} between them is far cheaper than a
process switch---no page tables to swap. But the kernel can schedule
threads at any instruction boundary, so two threads can interleave
their operations arbitrarily. That interleaving is the concurrency
problem, and only disciplined coordination makes it safe.
\end{osconcept}

\section{pthread\_create and pthread\_join}

The POSIX threads library (\code{pthreads}) is the C standard for
threads. A thread is created with \code{pthread\_create}, which starts
a function in a new thread, and \emph{reaped} with
\code{pthread\_join}, which waits for that thread to finish and
collects its result.

\begin{lstlisting}[caption={Creating two threads and waiting for them.}]
#include <stdio.h>
#include <stdlib.h>
#include <pthread.h>

/* A thread runs one function to completion. */
void *work(void *arg) {
    int id;

    id = *((int *) arg);
    printf("worker %d on thread %lu\n", id, pthread_self());
    return NULL;
}

int main(void) {
    pthread_t t1;
    pthread_t t2;
    int id1;
    int id2;
    int r;

    id1 = 1;
    id2 = 2;

    r = pthread_create(&t1, NULL, work, &id1);
    if (r != 0) {
        fprintf(stderr, "pthread_create: %d\n", r);
        return 1;
    }
    r = pthread_create(&t2, NULL, work, &id2);
    if (r != 0) {
        fprintf(stderr, "pthread_create: %d\n", r);
        return 1;
    }

    /* The main thread waits for both to finish. */
    pthread_join(t1, NULL);
    pthread_join(t2, NULL);

    printf("both workers are done\n");
    return 0;
}
\end{lstlisting}

The signature is the first thing to master:

\begin{verbatim}
   pthread_create( &thread_id,     /* where to store the new thread  */
                   NULL,           /* attributes: NULL = defaults    */
                   work,           /* the function to run            */
                   &id1 );         /* ONE pointer-sized argument     */

   - The function is void *work(void *arg): it takes ONE void* and
     returns ONE void*. That single pointer carries all input and
     output. This is the same "one value in, one value out" contract
     as thread arguments everywhere.

   - pthread_create returns an ERROR NUMBER, not -1, and it does NOT
     set errno. Check it against 0 and print the number itself.

   - pthread_join(thread, &result) blocks until that thread returns;
     the thread's return value is written through the second argument.

   - The argument we pass must OUTLIVE the thread. We pass the
     address of id1, a local of main that is still alive when the
     thread runs. Passing the address of a loop variable that
     changes would be a bug -- every thread would see the last value.
\end{verbatim}

\begin{verbatim}
   The lifetime contract, drawn:

     main
       |
       +-- pthread_create(&t1, work, &id1)  ->  thread t1 starts
       |                                        running work(id1)
       +-- pthread_create(&t2, work, &id2)  ->  thread t2 starts
       |                                        running work(id2)
       |
       +-- pthread_join(t1)  <-- main blocks here until t1 returns
       +-- pthread_join(t2)  <-- ...and here until t2 returns

   A thread is like a process child, but inside the same memory.
   create = fork, join = waitpid -- with sharing instead of copying.
\end{verbatim}

\begin{osconcept}
Thread creation and reaping are, again, OS services: the kernel must
allocate a new stack, new register set, and a new scheduler entry for
the thread, then run it on whatever CPU is free. The shared address
space is what makes threads cheap and also what makes them dangerous:
the kernel gives threads \emph{no isolation}, so the safety that
processes had by construction must now come from your discipline.
\end{osconcept}

\section{Data Races: The Shared-Variable Problem}

Here is the problem that defines all of concurrency. Two threads read
and write the same variable with no coordination. The C standard calls
this a \emph{data race}, and a data race is \emph{undefined behavior}--
the compiler may assume it never happens, and can optimize your
program in ways that make the bug appear far away from its cause.

\begin{lstlisting}[caption={The classic race: two threads, one counter.}]
#include <stdio.h>
#include <pthread.h>

#define N 1000000L

static long counter = 0;

void *increment(void *arg) {
    long i;

    for (i = 0; i < N; i++) {
        counter = counter + 1;
    }
    return NULL;
}

int main(void) {
    pthread_t t1;
    pthread_t t2;

    pthread_create(&t1, NULL, increment, NULL);
    pthread_create(&t2, NULL, increment, NULL);
    pthread_join(t1, NULL);
    pthread_join(t2, NULL);

    printf("counter = %ld (expected %ld)\n", counter, 2 * N);
    return 0;
}
\end{lstlisting}

Every run prints a different number, almost never \code{2000000}.
Why? Because \code{counter = counter + 1} is not one operation. It is
three: read, add, write. The kernel can interrupt a thread between any
two of them.

\begin{verbatim}
   The three steps of "counter = counter + 1", interleaved badly:

   Thread 1                    Thread 2
   -------                     -------
   read counter  (0)
   compute 0+1   (1)
   write counter (1)
                               read counter  (0)   <-- STALE: sees 0
                               compute 0+1   (1)
                               write counter (1)

   Both read 0, both write 1. One increment was LOST.
   With a million increments per thread, thousands are lost.

   A DATA RACE: two threads access the same memory, at least one
   write, with no synchronization in between. Result: undefined
   behavior. The number can be anything.
\end{verbatim}

The fix is not ``make the threads less unlucky.'' The fix is to make
the read-modify-write a single, indivisible unit---what we call a
\emph{critical section}---protected by a lock. That is the next
section.

\begin{ruleofthumb}[Any shared, written memory needs a lock]
The moment two threads can read a memory location that another thread
writes, you have a data race, and you have undefined behavior. The
only fix is synchronization: a mutex, a condition variable, or an
atomic operation. "It probably works" is not a fix; a race is a time
bomb with a random fuse.
\end{ruleofthumb}

\section{Mutexes: Making a Section Critical}

A \emph{mutex} (from ``mutual exclusion'') is a lock. Before entering
a critical section, a thread calls \code{pthread\_mutex\_lock}; while
one thread holds the lock, any other thread trying to lock is
\emph{blocked} by the kernel until the owner unlocks. The lock turns
``read, add, write'' into a section that only one thread may be inside
at a time.

\begin{lstlisting}[caption={The race, fixed with a mutex.}]
#include <stdio.h>
#include <pthread.h>

#define N 1000000L

static long counter = 0;
static pthread_mutex_t lock = PTHREAD_MUTEX_INITIALIZER;

void *increment(void *arg) {
    long i;

    for (i = 0; i < N; i++) {
        pthread_mutex_lock(&lock);       /* enter critical section */
        counter = counter + 1;           /* now safe: exclusive    */
        pthread_mutex_unlock(&lock);     /* leave critical section */
    }
    return NULL;
}

int main(void) {
    pthread_t t1;
    pthread_t t2;

    pthread_create(&t1, NULL, increment, NULL);
    pthread_create(&t2, NULL, increment, NULL);
    pthread_join(t1, NULL);
    pthread_join(t2, NULL);

    printf("counter = %ld (expected %ld)\n", counter, 2 * N);
    return 0;
}
\end{lstlisting}

\begin{verbatim}
   The critical section, drawn:

   Thread 1                          Thread 2
   lock()   -> owns the mutex
   counter++                         lock()   -> BLOCKS (kernel parks it)
   unlock() -> releases
                                     lock()   -> now owns it
                                     counter++
                                     unlock()

   The lock makes the read-modify-write atomic WITH RESPECT TO OTHER
   THREADS USING THE SAME LOCK. Two threads using DIFFERENT locks to
   protect the SAME variable are still racing -- the lock only works
   if everyone agrees on it.
\end{verbatim}

The discipline has four parts, and each is a rule:

\begin{itemize}
  \item \emph{One lock per unit of shared state.} The lock is a
    contract: every access to \code{counter} goes through it, or the
    contract is broken.
  \item \emph{Lock and unlock always pair.} Every path that locks must
    unlock---including error paths. (This is where C's lack of
    ``finally'' makes discipline essential.)
  \item \emph{Hold the lock as briefly as possible.} The critical
    section should contain the minimum work, so threads wait the
    minimum time.
  \item \emph{Never block while holding a lock if you can avoid it.}
    Holding a lock while waiting for I/O can stall every other thread
    that wants the lock.
\end{itemize}

\begin{osconcept}
A mutex is a synchronization primitive whose guarantee comes from the
kernel. On Linux a \code{pthread\_mutex\_t} lives in user space, and
the uncontended lock and unlock path never touches the kernel; only
when a thread must actually wait does the kernel (through a futex) park
it instead of letting it spin. Either way the OS guarantees that
\emph{at most one} thread holds a given mutex, and that guarantee is
what turns a sequence of instructions into a section others cannot
enter.
\end{osconcept}

\section{Condition Variables: Waiting for a Condition}

A mutex protects a variable; a \emph{condition variable} lets a thread
\emph{wait for a state of the world}---``buffer is not empty,''
``worker is done''---and be woken when that state arrives. The
condition variable is always used \emph{together with} a mutex, and
its central primitive does something that looks contradictory:

\code{pthread\_cond\_wait(cond, mutex)} \emph{atomically releases the
mutex and sleeps}. Another thread, seeing the state change, calls
\code{pthread\_cond\_signal} to wake one waiter. The awakened thread
re-acquires the mutex and re-checks the condition.

The canonical example is the bounded buffer: a fixed-size ring of
slots shared by producers and consumers.

\begin{lstlisting}[caption={Producer/consumer with a mutex and condition variables.}]
#include <stdio.h>
#include <pthread.h>

#define CAPACITY 8

static pthread_mutex_t mutex = PTHREAD_MUTEX_INITIALIZER;
static pthread_cond_t  not_full = PTHREAD_COND_INITIALIZER;
static pthread_cond_t  not_empty = PTHREAD_COND_INITIALIZER;

static int buffer[CAPACITY];
static int count = 0;

/* Produce one value into the shared buffer. */
void producer_put(int value) {
    pthread_mutex_lock(&mutex);

    /* Full? Then sleep until a consumer frees a slot.
       The loop, not an if, is mandatory: cond_wait can wake
       spuriously, and the state can change again before we
       actually run. */
    while (count == CAPACITY) {
        pthread_cond_wait(&not_full, &mutex);
    }

    buffer[count] = value;
    count = count + 1;

    pthread_cond_signal(&not_empty);   /* wake one consumer */
    pthread_mutex_unlock(&mutex);
}

/* Consume one value from the shared buffer. */
int consumer_get(void) {
    int value;

    pthread_mutex_lock(&mutex);

    /* Empty? Then sleep until a producer adds a value. */
    while (count == 0) {
        pthread_cond_wait(&not_empty, &mutex);
    }

    count = count - 1;
    value = buffer[count];

    pthread_cond_signal(&not_full);    /* wake one producer */
    pthread_mutex_unlock(&mutex);

    return value;
}
\end{lstlisting}

\begin{verbatim}
   The wait dance, step by step (producer side):

     producer locks the mutex
        |
        |  buffer full?
        |     yes --> cond_wait(not_full, mutex):
        |              1. atomically UNLOCKS the mutex
        |              2. sleeps until signaled
        |              3. on wake, RE-LOCKS the mutex
        |                     (then the while re-checks: full again?
        |                      wait again -- this is the loop's job)
        v
     slot is free: write, signal, unlock.

   The atomic unlock-and-sleep in cond_wait is the whole trick. If it
   released the lock and THEN slept, a producer could slip in between
   and the wait would race. The kernel makes the two a single step.
\end{verbatim}

\begin{ruleofthumb}[Always wait in a while loop, never an if]
\code{pthread\_cond\_wait} may return \emph{spuriously} (the kernel
does not guarantee one signal per wakeup), and by the time the waking
thread runs, the condition may have changed again. The correct shape
is always \code{while (condition) cond\_wait(...)}---re-check, don't
assume. This one habit eliminates an entire class of concurrency bugs.
\end{ruleofthumb}

\begin{osconcept}
A condition variable is the kernel's ``park me until the world
changes'' primitive, and its atomic release-and-sleep is the 
linchpin. Only the kernel can guarantee that the mutex is released
and the thread put to sleep as one indivisible step, so that no
signal is lost between the two. This is textbook operating-system
synchronization---semaphores and monitors are the same idea under
other names---and it is exactly why condition variables are kernel
primitives rather than library conveniences.
\end{osconcept}

\section{Deadlock: The Lock You Will Never Get}

When two threads each hold a lock the other needs, neither can
proceed, and both wait forever. That is \emph{deadlock}. It is the
second great concurrency bug, and it does not show up in tests---it
shows up in production, at 3 a.m., when two requests hit the locks in
the wrong order.

\begin{verbatim}
   Thread 1                      Thread 2
   -------                       -------
   lock(A)   ok                  lock(B)   ok
   lock(B)   BLOCKS              lock(A)   BLOCKS
             (B is held           (A is held
              by T2)               by T1)
   ...forever...                 ...forever...

   The deadly embrace, drawn as a cycle:

       T1 holds A, wants B
       T2 holds B, wants A
       --> a cycle of waiters. No one can make progress.

   The cure is ORDER: give every lock a global rank, and acquire
   locks in ascending order only. If everyone locks A before B,
   the cycle above cannot form.
\end{verbatim}

The four necessary conditions for deadlock are worth memorizing,
because removing \emph{any one} of them prevents it:

\begin{verbatim}
   deadlock needs all four:
     1. mutual exclusion     -- the locks exist
     2. hold and wait        -- a thread holds one lock, wants another
     3. no preemption        -- you cannot steal a held lock
     4. circular wait        -- the waiters form a cycle

   We cannot remove 1 (locks must exclude). We control 2 and 4 by
   design: acquire all locks in a fixed order (kills the cycle), and
   hold nothing while waiting for I/O (kills hold-and-wait).
\end{verbatim}

\begin{lstlisting}[caption={The deadlock-free pattern: lock in one order, always.}]
#include <pthread.h>

void transfer(pthread_mutex_t *from_lock, pthread_mutex_t *to_lock,
              long *from, long *to, long amount) {
    /* Both threads lock in the SAME order (lowest address first).
       Because the order is fixed, no cycle of waiters can form. */
    if (from_lock < to_lock) {
        pthread_mutex_lock(from_lock);
        pthread_mutex_lock(to_lock);
    } else {
        pthread_mutex_lock(to_lock);
        pthread_mutex_lock(from_lock);
    }

    *from = *from - amount;
    *to = *to + amount;

    pthread_mutex_unlock(to_lock);
    pthread_mutex_unlock(from_lock);
}
\end{lstlisting}

Comparing lock \emph{addresses} is a real, used technique: it imposes
a total order on locks without naming them. Whatever the two threads
pass in, they acquire them in the same sequence, and the circular wait
becomes impossible.

\begin{osconcept}
Deadlock is the kernel being \emph{too} faithful: it parks a thread
forever because that is exactly what the lock primitive promises. The
OS cannot detect and fix your lock ordering for you---it only enforces
mutual exclusion. Recognizing the four conditions and designing them
away is a programming discipline, not a runtime feature, which is why
deadlock prevention appears in operating-system textbooks and in
every serious concurrency course.
\end{osconcept}

\section{Threads vs Processes: Which One}

With both toolkits in hand, the choice is a design decision, not a
religious one. The table is the decision procedure:

\begin{verbatim}
   concern            threads                      processes
   ----------------   ---------------------------  -------------------
   address space      ONE, shared                  private each
   sharing data       variables (with locks)       pipes, shared memory
   creation cost      cheap (a stack + TCB)        heavier (page tables)
   crash isolation    NONE: one thread's segfault  strong: siblings
                      kills the whole process      survive
   communication      implicit, via memory         explicit, via kernel
   best for           parallel compute, latency-   robustness, multi-
                      critical I/O workers, GUI    tenant, "one bug
                                                    must not spread"

   The rule of thumb:
     Need to move lots of data between two workers? THREADS.
     Need one worker to be immune to another's crash? PROCESSES.
     Need both? The classic pattern: a process per service, threads
     inside each service.
\end{verbatim}

\section{Practical: A Parallel Sum}

Let us put the chapter to work with a real parallel algorithm: sum a
million numbers with four threads, each summing its own slice, then
combine. The trick is that each thread writes only to its \emph{own}
range struct---no shared counters, no locks needed, because the
sharing is done once, up front, by careful partition.

\begin{lstlisting}[caption={parsum.c: four threads, four disjoint slices.}]
#include <stdio.h>
#include <stdlib.h>
#include <pthread.h>

#define NTHREADS 4
#define ELEMENTS 1000000

static int data[ELEMENTS];

typedef struct range {
    int start;          /* first index, inclusive  */
    int end;            /* last index, exclusive   */
    long total;         /* this slice's sum        */
} range;

void *sum_range(void *arg) {
    range *r;
    long acc;
    int i;

    r = (range *) arg;
    acc = 0;
    for (i = r->start; i < r->end; i++) {
        acc = acc + data[i];
    }
    r->total = acc;          /* only this thread touches r->total */
    return NULL;
}

int main(void) {
    pthread_t threads[NTHREADS];
    range ranges[NTHREADS];
    long grand;
    size_t i;

    /* Fill the array. */
    for (i = 0; i < (size_t) ELEMENTS; i++) {
        data[i] = 1;
    }

    /* Split the array into NTHREADS disjoint slices. */
    for (i = 0; i < (size_t) NTHREADS; i++) {
        ranges[i].start = (int) (i * ((size_t) ELEMENTS / NTHREADS));
        ranges[i].end = (int) ((i + 1) * ((size_t) ELEMENTS / NTHREADS));
        ranges[i].total = 0;
        pthread_create(&threads[i], NULL, sum_range, &ranges[i]);
    }

    /* Collect: join each thread, add its slice. */
    grand = 0;
    for (i = 0; i < (size_t) NTHREADS; i++) {
        pthread_join(threads[i], NULL);
        grand = grand + ranges[i].total;
    }

    printf("sum = %ld (expected %d)\n", grand, ELEMENTS);
    return 0;
}
\end{lstlisting}

\begin{verbatim}
   Why this program is race-free WITHOUT a single lock:

     thread 0  owns data[0 .. 249999]   and ranges[0].total
     thread 1  owns data[250000 .. 499999]  and ranges[1].total
     ...
     the main thread writes ranges[].start/.end BEFORE creating
     the threads, and reads ranges[].total AFTER joining them.

   Every memory location is written by exactly ONE thread at a time.
   The only synchronization needed is the create/join itself, which
   pthreads provides. Shared nothing until the combine--the design
   discipline that makes concurrency easy.
\end{verbatim}

\begin{researchnote}[The origin of mutual exclusion]
The mutex and condition variable you just used descend directly from
work done in 1965 by Edsger Dijkstra, in a technical report called
``Cooperating Sequential Processes.'' In it, Dijkstra posed the
\emph{mutual exclusion problem}---two processes sharing variables
must not interleave their critical sections---and introduced the
\emph{semaphore}, the ancestor of the mutex, along with the very idea
of a critical section.

\begin{verbatim}
   Dijkstra's formulation, condensed:

       "If two processes share variables, each must protect its
        critical section so that at most one is inside at any time;
        the semaphore is the primitive that guarantees it."

   Our pthread_mutex_lock/unlock and cond_wait are the modern,
   clarified descendants of those first semaphores. When you write a
   critical section, you are using a concept that is sixty years old
   and still state of the art.
\end{verbatim}

It is a fascinating document---early concurrency worked out by hand,
before anyone had seen it fail in production. Reading it is the
fastest way to understand why the rules of this chapter exist.

\paperref{E.~W.~Dijkstra, ``Cooperating Sequential Processes,''
Technical Report EWD-123, Technological University Eindhoven,
1965}{https://www.cs.utexas.edu/users/EWD/transcriptions/EWD01xx/EWD123.html}
\end{researchnote}

\section{Structures in Bytes}

The pthread objects are user-space byte structures whose exact layout
is an ABI detail you should treat as opaque---but their shapes are
instructive. On glibc/x86-64:

\begin{verbatim}
   pthread_t       8 bytes    an opaque handle
   pthread_mutex_t 40 bytes   owner + kind + count + wait list
   pthread_cond_t  48 bytes   wait list + clock + sequencing

   pthread_mutex_t (glibc x86-64), approximate offsets:
   offset  size  field
   0       4     lock
   4       4     count        (recursion depth)
   8       4     owner        (which thread holds it)
   12      4     nusers
   16      4     kind         (normal, recursive, errorcheck)
   20      2     spins
   24      8     wait list    (futex waiters)
\end{verbatim}

When the kernel switches threads, the byte structure it saves is the
register image: 15 general-purpose registers at 8 bytes each (120
bytes) plus the instruction and stack pointers, stored on the thread's
kernel stack. That saved image, swapped on every context switch, is the
mechanical meaning of ``two threads, two sets of registers, one
address space.''

\section{Summary}

\begin{verbatim}
   The concurrency toolbox, complete:

   pthread_create / pthread_join   spawn and reap threads
   shared address space            threads share heap and globals
   data race                       unsynchronized shared writes = UB
   pthread_mutex_t                 the critical section (mutual excl.)
   pthread_cond_t                  wait for a condition, with the lock
   deadlock                        lock cycles; fix by lock ordering

   Rules that keep threaded code correct:
     1. every shared, written variable is protected by a lock;
     2. lock and unlock always pair, on every path;
     3. hold locks briefly; never block for I/O while holding one;
     4. cond_wait in a while loop, never an if;
     5. acquire multiple locks in one global order;
     6. prefer disjoint ownership (like the parallel sum) so locks
        are rarely needed at all.
\end{verbatim}

\begin{osconcept}
Threads are the OS's cheapest form of parallelism, purchased by
surrendering the isolation that processes enjoyed. The kernel
multiplexes instruction streams and enforces the primitives (locks,
condvars) that make sharing safe; everything else---partitioning,
ordering, and the discipline of the critical section---is the
programmer's. Master that bargain and you can write programs that use
every core; break it and you get bugs that reproduce only in
production.
\end{osconcept}

The next chapter takes the sharing bargain to its logical extreme:
shared memory, where processes---not threads---map the same bytes into
their address spaces and synchronize over raw memory, with no kernel
copying and no safety net.

